\begin{tabular}{p{0.20\textwidth}p{0.36\textwidth}p{0.36\textwidth}}
\toprule
Блок & Основной количественный результат & Граница применимости \\
\midrule
Сопоставление архитектур & Обучаемое правило по оцененному состоянию улучшает результат относительно жестко заданного правила на 56.7--71.1\% по четырем сценариям; прямая политика по наблюдениям улучшает его на 24.7--73.6\%. & Прямая политика по наблюдениям полезнее только в сценариях с более богатым набором информации; при тонком информационном наборе лучше работает обучаемое правило по отфильтрованному состоянию. \\
Карта архитектурных ошибок & Лучшее обучаемое правило превосходит классические схемы с ошибками фильтрации и задания вероятностей режима в среднем на 62.9--66.6\%; доля выигрыша по сценариям и по проверочным траекториям равна единице. & Преимущество не является универсальным: против удачного простого правила, реагирующего только на инфляцию, обучаемое правило уже не выигрывает. \\
Перенос на новые среды & При переносе на новые среды обучаемое правило остается лучше фиксированного правила в среднем на 51.9--61.6\%; выигрыш наблюдается во всех средах. & Однако перенастроенное простое правило переносится лучше: обучаемое правило проигрывает ему во всех средах и теряет больше качества при переносе. \\
\bottomrule
\end{tabular}